\documentclass[language=ngerman,algolist,acrolist]{TemplateLecture}

\renewcommand{\LectureName}{Grundlagen der algorithmischen Geometrie}
\renewcommand{\ProfName}{Dr. Herman Haverkort, \\ PD. Dr. Elmar Langetepe}
\renewcommand{\ProfNameShort}{Haverkort, Langetepe}
\renewcommand{\Semester}{Sommer 2026}
\renewcommand{\mName}{Daniel Emse}


\acsetup{
  index/use = true
}

\DeclareAcronym{es}{
    short=ES,
    long=Ereignisstruktur,
}

\DeclareAcronym{sss}{
    short=SSS,
    long=Sweep-Status-Struktur,
}

\DeclareAcronym{avl}{
    short=AVL,
    long={Adelson-Velsky, Landis}
}










\newcommand{\MaxSF}{\textit{MaxSoFar}}
\newcommand{\MinSF}{\textit{MinSoFar}}
\newcommand{\links}{\textit{links}}
\newcommand{\rechts}{\textit{rechts}}
\newcommand{\MinDistSF}{\textit{MinDistSoFar}}
\newcommand{\MaxEndHere}{\textit{MaxEndingHere}}
\newcommand{\ClosPos}{\textit{ClosPos}}
\newcommand{\event}{\textit{Ereignis}}
\newcommand{\MinDist}{\textsc{MinDist}}
\newcommand{\Insert}{\textsc{Insert}}
\newcommand{\Delete}{\textsc{Delete}}
\newcommand{\Prec}{\textsc{Prec}}
\newcommand{\Succ}{\textsc{Succ}}
\newcommand{\NextEvent}{\textsc{NächstesEreignis}}

\usepackage{algpseudocodex}

\begin{document}
\setcounter{chapter}{1}

%%%%%%%%%%%%%%%%%%%%%%%%%%%%%%%%%%%%%%%%%%%%%%%%%%%%%%%%%%%%%%%%%%%%%%%%%%%%%%%%%%%%%%%%%%%%%%%%%%%%
%%%%%%%%%%%%%%%%%%%%%%%%%%%%%%%%%%%%%%%%%% 2026 / 04 / 14 %%%%%%%%%%%%%%%%%%%%%%%%%%%%%%%%%%%%%%%%%%
%%%%%%%%%%%%%%%%%%%%%%%%%%%%%%%%%%%%%%%%%%%%%%%%%%%%%%%%%%%%%%%%%%%%%%%%%%%%%%%%%%%%%%%%%%%%%%%%%%%%

\chapter{Das Sweep-Verfahren}
\section{Einführung}

\newLecture{14.04.2026}

Das \emph{Sweep}\index{Sweep}-Verfahren ist ähnlich vielseitig wie \emph{Divide and Conquer}, also ein hilfreiches Pardigma. Divide and Conquer ist beispielsweise sinnvoll, wenn die erwarteten Kosten des Ergebnisses klein sind. Beispielsweise erhalten wir $T(n) \leq 2 \cdot T(n/2) + c(\text{merge})$. 

\emph{Sweep} ist ein geometrisches Verfahren, bei dem wir lokal eine Dimensionsreduktion von einem $n$-dimensionalen Problem auf ein $(n-1)$-dimensionales Problem verwandeln, indem wir die Ebene mit einer \emph{Sweepline} ausfegen, indem wir gemäß einer vorher angefertigten Sortierung fegen. Dabei folgt die Korrektheit des Verfahrens über eine Invariante nahe der Sweepline.

\section{Sweep im Eindimensionalen}
\subsection{Das Maximum einer Menge von Objekten}

\begin{prob}{Find maximum}{findMax}
    \prblm{$Q = [q_1, \ldots, q_n] \in \IR^n$}{$q_i \in Q$ mit $q_i \geq q_j \forall j = 1,\ldots,n$}
\end{prob}

Hierfür haben wir folgenden optimalen Algorithmus mit Laufzeit $\Theta(n)$. Dabei sweepen wir über die Indizes, erhalten also lokal ein $0$-dimensionales Problem

\begin{algorithm}[h]
    \caption{\textsc{FindMax}}
    \begin{algorithmic}[1]
        \Input{Wie in~\Cref{prob:findMax}}
        \Output{Wie in~\Cref{prob:findMax}}
        \State{$\MaxSF \gets q[1]$}
        \For{$j = 2$ \textbf{to} $n$ \textbf{do}}
            \If{$\MaxSF < q[j]$}
                \State{$\MaxSF \gets q[j]$}
            \EndIf{}
        \EndFor{}
        \State{\Return{$\MaxSF$}}
    \end{algorithmic}
\end{algorithm}

\subsection{Das dichteste Paar in einer Menge von Zahlen}

 
\begin{prob}{Closest pair in einer Dimension}{closestPair1D}
    \prblm{$x_1, \ldots, x_n \in \IR$}{$i \neq j \in \set{1,\ldots,n}$ sodass $\displaystyle\abs{x_i - x_j} = \min_{\ell \neq k} \abs{x_\ell - x_k}$ für $\ell, k \in \set{1,\ldots,n}$}
\end{prob}

Wir sortieren die Zahlen aufsteigend, sodass wir ein neues Array $x_1' \leq  \cdots \leq x_n' \in \IR$ bekommen (im Mehrdimensionalen sortieren wir sie nach einer bestimmten Koordinate). Das braucht aber natürlich bereits $\cO(n \log(n))$ Zeit.

\begin{algorithm}[h]
    \caption{\textsc{FindClosestPair-1D}}
    \begin{algorithmic}[1]
        \Input{Wie in~\Cref{prob:closestPair1D}, aber als aufsteigend sortiertes Array $x'$}
        \Output{Wie in~\Cref{prob:closestPair1D}}
        \State{$\MinDistSF \gets x'[2] - x'[1]$}
        \State{$\ClosPos \gets 2$}
        \For{$j = 3$ \textbf{to} $n$ \textbf{do}}
            \If{$\MinDistSF > x'[j] - x'[j-1]$}
                \State{$\MinDistSF \gets x'[j] - x'[j-1]$}
                \State{$\ClosPos \gets j$}
            \EndIf{}
        \EndFor{}
        \State{\Return{$x'[\ClosPos-1], x'[\ClosPos]$}}
    \end{algorithmic}
\end{algorithm}

\begin{cor}
    Die Berechnung des \textsc{Closest Pair}s von $n$ reellen Zahlen benötigt Laufzeit $\cO(n \log(n))$.
\end{cor}

\subsection{Die maximale Teilsumme}

\begin{prob}{Maximum subvector}{maximumSubvector}
    \prblm{Array $V \in \IR^n$}{Finde $i \leq j$, sodass $\displaystyle V[i] + V[i+1] + \cdots + V[j] = \max_{\ell \leq k} V[\ell] + \cdots + V[k]$, mit $i, j, \ell, k \in \set{1,\ldots,n}$}
\end{prob}

\begin{algorithm}[h]
    \caption{\textsc{FindMaximumSubvector}}
    \begin{algorithmic}[1]
        \State{$\MaxSF \gets 0$}
        \State{$\MaxEndHere \gets 0$}
        \For{$j = 1$ \textbf{to} $n$ \textbf{do}}
            \State{$\MaxEndHere \gets \max(0, \MaxEndHere + V[j])$}
            \State{$\MaxSF \gets \max(\MaxSF, \MaxEndHere)$}
        \EndFor{}
        \State{\Return{$\MaxSF$}}
    \end{algorithmic}
\end{algorithm}

\begin{thm}{FindMaximumSubvector}{findMaxSubvec}
    In einer Folge von $n$ Zahlen die maximale Teilsumme konsekutiver Zahlen zu bestimmen, hat die Zeitkomplexität $\Theta(n)$.
\end{thm}



\section{Sweep in der Ebene}

Im Allgemeinen ist $\abs{x} \coloneqq \norm{x}_2$ die Euklidische Norm des Vektor $x$.

\subsection{Das dichteste Punktepaar in der Ebene}
 
\begin{prob}{Closest pair in zwei Dimensionen}{closestPair2D}
    \prblm{$x_1, \ldots, x_n \in \IR^2$}{$i \neq j \in \set{1,\ldots,n}$ sodass $\displaystyle\abs{x_i - x_j} = \min_{\ell \neq k} \abs{x_\ell - x_k}$ für $\ell, k \in \set{1,\ldots,n}$}
\end{prob}

Wir führen zum Verwalten ein Verzeichnis ein, in dem zu jedem Zeitpunkt den Status der \emph{Sweep Line} tracked. Dieses heißt \emph{\ac{sss}}. Übergeordnet führen wir eine \emph{\ac{es}} ein, in dieser sind einfach alle Punkte nach ihrer $x$-Koordinate sortiert. Die \ac{sss} teilt diese in Streifen der Breite $\MinSF$ ein. Erreicht ein neuer Punkt solch einen Streifen, testen wir unter allen Punkten im Streifen, welche davon den geringsten Abstand haben. Das aktualisiert $\MinSF$ und der Streifen der \ac{sss} wird entsprechend schmaler.

Wir erhalten den folgenden Algorithmus:

\begin{algorithm}[h]
    \caption{\textsc{FindClosestPair-2D}}
    \begin{algorithmic}[1]
        \Input{Array von $n$ Punkten aus $\IR^2$}
        \Output{Zwei Elemente des Arrays mit dem geringsten euklidischen Abstand}
        \State{Sortiere die $n$ Punkte nach aufsteigenden $x$-Koordinaten und füge sie ins Array $P$ ein. Das bildet die \ac{es}}
        \State{Füge $P[1], P[2]$ in \ac{sss} ein}
        \State{$\MinSF \gets \abs{P[1] - P[2]}$}
        \State{$\links \gets 1$}
        \State{$\rechts \gets 3$}
        \While{$\rechts \leq n$}
            \If{$P[\links].x + \MinSF \leq P[\rechts].x$}
                \State{Entferne $P[\links]$ aus \ac{sss}}
                \State{$\links \gets \links + 1$}
            \Else{}
                \State{$\MinSF \gets \MinDist(\text{\ac{sss}}, P[\rechts], \MinSF)$}
                \State{Füge $P[\rechts]$ in \ac{sss} ein}
                \State{$\rechts \gets \rechts + 1$}
            \EndIf{}
        \EndWhile{}
        \State{\Return{$\MinSF$}}
    \end{algorithmic}
\end{algorithm}


\newLecture{16.04.2026}

Zu spezifizieren ist noch \MinDist. Wir müssen uns dafür Gedanken darüber machen, welche Datenstruktur wir am besten verwenden, um die \ac{sss} zu verwalten. Eine Liste hätte lineare Laufzeit, um die Nachbarn zu prüfen, das kann im schlimmsten Fall $\cO(n^2)$ pro Ausführung von \MinDist benötigen kann.

Sei $r \coloneqq P[\rechts]$ und $M \coloneqq \MinSF$. Dann genügt es, für $\MinDist(\text{\ac{sss}}, r, M)$ nur die Punkte aus \ac{sss} anzuschauen, deren $x$-Koordinate in $\edgys{r.x - M, r.x + M}$ liegt. Folgende Operationen muss unsere \ac{sss} unterstützen:
\begin{itemize}
    \item $\MinDist(\text{\ac{sss}}, r, M)$
    \item $\Insert(p)$
    \item $\Delete(p)$
\end{itemize}
Wir verwenden als Datenstruktur einen balancierten Binärbaum über die $y$-Koordinaten als Schlüssel. Damit benötigen $\Insert$ und $\Delete$ jeweils $\cO(\log(n))$, und Ausführung der Bereichsanfrage nach allen $p.y$ im Invervall $\edgys{r.x - M, r.x + M}$ benötigt Laufzeit $\cO(\log(n) + k)$, wobei $k$ die Größe der Antwort bedeutet.

Der betrachtete Streifen des Rechtecks um $r$ hat einen Flächeninhalt von $2M^2 = 2M \cdot M$. Ein Punkt aus dem Rechteck deckt mit seinem Kreis mit Radius $\frac{M}{2}$ mindestens einen Viertelkreis als Flächeninhalt des Rechtecks ab, also mindestens $\pi\frac{M^2}{16}$. Wir erhalten mit
\begin{equation*}
    k \leq \frac{2M^2}{\pi\frac{M^2}{16}} \leq \frac{32}{\pi} < 11,
\end{equation*}
also ist $k \leq 10\in \cO(1)$ und die Bereichsanfrage benötigt auch nur $\cO(\log(n))$ Zeit. Die Datenstruktur für \ac{sss} in diesem Fall ist also ein \emph{\ac{avl}-Baum}. Wir erhalten folgendes Lemma:

\begin{lem}{Höchstens $10$ Punkte in Rechteck}{max10ptsInRect}
    Sei $M > 0$ und $P$ eine Menge von Pnukten im $\IR^2$, von denen je zwei den Abstand $\geq M$ haben. Ein Rechteck mit Kantenlängen $M$ und $2M$ enthält höchstens $10$ Punkte.
\end{lem}

\begin{thm}{Laufzeit FindClosestPair}{findClosestPairRuntime}
    Der minimale Abstand aller Paare einer $n$-elementigen Punktmenge in der Ebene lässt sich in Zeit $\cO(n \log(n))$ bestimmen.
\end{thm}

\subsection{Schnittpunkte von Strecken}

\begin{prob}{Line Segment Intersection in zwei Dimensionen}{lineSegmentIntersection2D}
    \prblm{$n$ Liniensegmente $s_i = l_i r_i$ in der Ebene für $i \in \set{1,\ldots,n}$}{Berechne die \emph{echten Schnittpunkte}, also Punkte, die zum relativen Inneren beider Strecken bestehen}
\end{prob}
Wir nehmen an, dass keine Mehrfachschnittpunkte vorliegen und nie gleiche $x$-Koordinaten an beiden Segmentenden vorkommen. Außerdem sollen die geraden in \emph{allgemeiner Lage}\index{allgemeine Lage} liegen, sie weisen also keinerlei geometrische Besonderheiten wie Parallelität oder Orthogonalität auf.

Im \emph{Existenzproblem}\index{Existenzproblem} fragen wir uns, ob es unter den $n$ Strecken mindestens zwei gibt, die einen echten Schittpunkt haben, bein \emph{Aufzählungsproblem}\index{Aufzählungsproblem} wollen wir \emph{alle} echten Schnittpunkte bestimmen.

Wir fragen uns wieder nach einer sinnvollen Wahl von \acl{es} und \acl{sss}: Die \ac{es} bestehe aus den linken und rechten Endpunkten $l_i$ und $r_i$, insbesondere deren $x$-Koordinaten. Als \ac{sss} wählen wir die Segmente gemäß einer noch festzulegenden Ordnung. Außerdem gibt es drei Typen von Ereignissen, die auftreten können:
\begin{enumerate}
    \item Neuer linker Endpunkt eines Segments gefunden (\LeftEnd),
    \item Neuer rechter Endpunkt eines Segments gefunden (\RightEnd),
    \item Neuer Schnittpunkt gefunden (\Intersect).
\end{enumerate}

\begin{lem}{Echte Schnittpunkte sind Nachbarn}{intersectAreNeighbours}
    Haben zwei Liniensegmente einen echten Schnittpunkt zum Zeitpunkt $t$, dann sind die Segmente unmittelbar vor $t$ Nachbarin in der \ac{sss} geworden.
\end{lem}

wir fragen uns jetzt, welche Operationen wir für die Lösung des \emph{Existenzproblems} unterstützen müssen:
\begin{itemize}
    \item $\Insert(\text{\ac{sss}}, seg, x)$ fügt Segment $seg$ entsprechend der Ordnung zur Zeit $x$ in die \ac{sss} ein.
    \item $\Delete(\text{\ac{sss}}, seg, x)$ entfernt Segment $seg$ zur Zeit $x$ aus der \ac{sss}.
    \item $\Prec(\text{\ac{sss}}, seg, x)$ bestimmt den Vorgänger von $seg$ in der \ac{sss} zum Zeitpunkt $x$.
    \item $\Succ(\text{\ac{sss}}, seg, x)$ bestimmt den Nachfolger von $seg$ in der \ac{sss} zum Zeitpunkt $x$.
    \item $\IntersectTest(seg_1, seg_2)$ testet auf Schnitt zwischen $seg_1$ und $seg_2$.
\end{itemize}

Für die Implementierun der \ac{sss} verwenden wir wieder einen \ac{avl}-Baum als balancierten Suchbaum, also ist jede der Operationen in $\cO(\log(n))$ ausführbar, weil höchstens $n$ Segmente in der \ac{sss} gespeichert sind und weil der Suchbaum balanciert ist. Da die Bearbeitung von jedem der höchstens $2n$ vielen Ereignisse nur drei solche Operationen erfordert, erhalten wir folgende Aussage:

\begin{thm}{Existenzproblem Line Segment Intersection}{existanceProblemLineSegmentIntersection}
    Man kann in optimaler Zeit $\cO(n \log(n))$ und mit Speicherplatz $\cO(n)$ herausfinden, ob von $n$ Strecken in der Ebene mindestens zwei einen echten Schnittpunkt haben.
\end{thm}

Dabei haben \Insert{} und \Delete{} wieder Laufzeit $\cO(\log(n))$, während $\IntersectTest$ nur $\cO(1)$ benötigt.

Das Lösen des \emph{Aufzählungsproblem} ist etwas anspruchsvoller. Nach einem Schnittpunkt wechselt die Ordnung zwischen zwei Segmenten, das vorher obere ist dann das neue untere Segment. Die konkreten Instanzen der drei oben beschriebenen Typen von Ereignissen, die während des Sweep zu behandeln sind, sind nicht alle schon von Anfang bekannt, deswegen brauchen wir eine \emph{dynamische} \ac{es}. Außer Initialisierung und dem Test, ob \ac{es} leer ist, sind folgende Operationen möglich:
\begin{itemize}
    \item $\Insert(\text{\ac{es}},\event)$ fügt das Ereignis \event{} entsprechend seiner Eintrittszeit in die \ac{es} ein.
    \item $\NextEvent$ liefert das zeitlich erste in \ac{es} gespeicherte Ereignis und entfernt es aus der \ac{es}.
\end{itemize}

Dafür werden wir eine \emph{Priority Queue} verwenden.

Zusammen haben wir folgenden Algorithmus:

\begin{algorithm}[h]
    \caption{\textsc{FindAllIntersections-2D}}\label{alg:findAllIntersections2D}
    \begin{algorithmic}[1]
        \State{Initialisiere \ac{sss} und \ac{es}}
        \State{Sortiere die $2n$ Endpunkte nach aufsteigenden $x$-Koordinaten}
        \State{Erzeuge daraus \event{}s und füge diese in die \ac{es} ein}
        \While{$\text{\ac{es}} \neq \emptyset$}
            \State{$\event \gets \NextEvent(\text{\ac{es}})$}
            \If{\event{} \textbf{is} \LeftEnd}
                \State{$\Insert(\text{\ac{sss}}, seg, t)$}
                \State{$prevSeg \gets \Prec(\text{\ac{sss}}, seg, t)$}
                \State{$\IntersectTestCreateEvent(prevSeg,seg)$}
                \State{$nextSeg \gets \Succ(\text{\ac{sss}}, seg, t)$}
                \State{$\IntersectTestCreateEvent(seg,nextSeg)$}
            \Else{} 
                \If{\event{} \textbf{is} \RightEnd}
                    \State{$prevSeg \gets \Prec(\text{\ac{sss}}, seg, t)$}
                    \State{$nextSeg \gets \Succ(\text{\ac{sss}}, seg, t)$}
                    \State{$\Delete(\text{\ac{sss}}, seg, t)$}
                    \State{$\IntersectTestCreateEvent(prevSeg,nextSeg)$}
                \Else{}
                    \State{$(uSeg, lSeg) \gets (\text{vorher oberes Segment}, \text{vorher unteres Segment})$}
                    \State{Nehme $(uSeg, lSeg)$ als Paar mit Schnitt auf}
                    \State{$prevSeg \gets \Prec(\text{\ac{sss}}, uSeg, t)$}
                    \State{$\IntersectTestCreateEvent(prevSeg,uSeg)$}
                    \State{$nextSeg \gets \Succ(\text{\ac{sss}}, uSeg, t)$}
                    \State{$\IntersectTestCreateEvent(lSeg,nextSeg)$}
                \EndIf{}
            \EndIf{}
        \EndWhile{}
    \end{algorithmic}
\end{algorithm}

\begin{thm}{FindAllIntersections-2D}{findAllIntersections2D}
    Das Sweep-Verfahren berechnet die $k$ echten Schnittpunkte von $n$ Liniensegmenten in Zeit $\cO((n + k) \log(n))$ mit Platz $\cO(k + n)$.
\end{thm}

\begin{proof}
    Insgesamt gibt es $2n+k$ Ereignisse. Weil keines von ihnen mehrfach in der Ereignis-Struktur vorkommt (da jeder Versuch, ein schon vorhandenes Element erneut einzufügen, von der Datenstruktur abgewiesen wird), sind niemals mehr als $\cO(n^2)$ Ereignisse in der \ac{es} gespeichert. Jeder Zugriff auf die \ac{es} ist also in Zeit $\cO(\log(n^2)) = \cO(\log(n))$ ausführbar.
    
    Die \textbf{while}-Schleife wird $(2n + k)$-mal durchlaufen; bei jedem Durchlauf werden höchstens sieben Operationen auf der \ac{sss} oder der \ac{es} ausgeführt. Daraus folgt die Behauptung.
\end{proof}

\begin{thm}{FindAllIntersections-2D weniger Platzbedarf}{findAllIntersections2DLessSpace}
    Das Sweep-Verfahren berechnet die $k$ echten Schnittpunkte von $n$ Liniensegmenten in Zeit $\cO((n + k) \log(n))$ mit Platz $\cO(k + n)$.
\end{thm}

Wir können also auch degenerierte Fälle behandeln.

\newLecture{21.04.2026}


\newLecture{23.04.2026}

\end{document}